\documentclass[main.tex]{subfiles}
\begin{document}
\label{introduction}

The COVID-19 pandemic raised fundamental questions about how to contain the spread
of a novel disease \cite{LepriCovid, EvaluationCovid} and about the collateral
effects of containment strategies \cite{CovidEffectEconomy}. In the absence of
vaccines, Non-Pharmaceutical Interventions (NPIs) were the primary tool adopted by
governments to limit transmission \cite{NPI1, NPIDD, NPI3}, significantly reshaping
human behaviour and mobility patterns \cite{SantanaChanges}. Importantly, a
substantial share of these behavioural changes preceded formal mandates, driven by
voluntary responses to risk information \cite{Goolsbee2020, Gupta2020, Xiong2020},
a distinction with direct consequences for how mobility-based models should be
calibrated. This context provided the scientific community with a large-scale natural
experiment for developing and evaluating models of epidemic spread
\cite{model0, model1, model2, model3, model4}.

Aggregate and anonymised mobile phone data have long been shown to support the
modelling of epidemic geography \cite{HumanProxiesModelling, MobilityMalaria}, while
individual GPS trajectories enable finer estimates of contact patterns and spatial
transmission dynamics \cite{Finger2016, LepriCovid, Alessandretti, Chang2021,
Kraemer2020}. The COVID-19 pandemic catalysed an unprecedented release of such data,
from aggregated public indices \cite{GoogleMobility, AppleMobility} to
privacy-enhanced individual-level datasets \cite{Chang2021, lucchini}, enabling
analyses at scales previously unavailable. These datasets are nonetheless subject to
demographic sampling biases that must be accounted for when drawing population-level
inferences \cite{Wesolowski2013, Zhao2016, Li2024}.

Mobility data are often partial, aggregated, or available at mismatched spatial
scales \cite{Gatto2020, Bertuzzo2020, Balcan2009a}, leading studies to complement
them with probabilistic mobility models \cite{Aleta2020, Arenas2020,
HumanProxiesModelling, CiofidegliAtti2008}. The validity of such models is
compromised when restriction measures alter the underlying mobility patterns
\cite{lucchini, Schlosser2020a}, and simple rescaling approaches prove inadequate.
The reduction in mobility was itself heterogeneous: sharp differences have been
documented along income gradients \cite{Weill2020}, urban--rural divides
\cite{Li2020}, and partisan lines \cite{Allcott2020, Grossman2020}. Recovery
dynamics after restrictions were lifted showed similar heterogeneity across
jurisdictions and political contexts \cite{Kim2021, Nguyen2020}. Understanding this
variability --- and whether standard statistical descriptions of human mobility
remain valid under forced restrictions --- is the central motivation of this work.

The relationship between trip counts, contacts, and transmission probability did not
remain constant across pandemic phases, driven by shifting levels of public awareness
\cite{nouvellet, rg_forecast_deaths}, making it challenging to infer viral spread
from mobility data alone \cite{Alessandretti}. We address this by providing a
comprehensive individual-level analysis across three dimensions: geometric measures,
visitation patterns, and information-theoretic complexity.

We analyse a large, privacy-enhanced longitudinal GPS dataset comprising over
$180{,}000$ anonymised trajectories from Massachusetts and over $1{,}000{,}000$ from
California, covering January to September 2020. The observation period is divided
into three phases \cite{lockdown}: pre-restriction (15 January -- 15 March),
lockdown (15 March -- 15 May, when more than 40 U.S.\ states issued shelter-in-place
orders), and post-restriction (15 May -- end of September).

We measure the radius of gyration \cite{radius_gyration} and test whether county-level
political affiliation and urbanisation modulate it across phases. Trip-length
distributions follow a power law \cite{radius_gyration, brockel} whose shape we
examine under restriction. The rate of discovery of new places follows a diffusive
power law \cite{scaling_visitation}, and the rank structure of visitation frequencies
obeys a universal law \cite{scaling_visitation, Schlapfer2021} which we test for
stability. We further compute stop-point distributions in each trajectory's intrinsic
reference frame \cite{radius_gyration}. For trajectory complexity, we use three
entropy measures --- random, Shannon, and Lempel-Ziv real entropy
\cite{entropic_measures, Lu2013} --- which progressively add sensitivity to visit
frequency, sequential structure, and the regularity of return behaviour
\cite{Teixeira2021, Cuttone2016}. Whether complexity reduction during lockdown
reflects a structural change in trajectories or simply fewer visited locations is one
of the central questions we address. Finally, we apply a returner--explorer framework
\cite{pappalardo_returners} to characterise shifts in habitual versus novel movement.

\end{document}